\documentclass[11pt,a4paper]{article}

\usepackage[margin=1in]{geometry}
\usepackage[T1]{fontenc}
\usepackage[utf8]{inputenc}
\usepackage{lmodern}
\usepackage{microtype}
\usepackage[round,authoryear]{natbib}
\usepackage[colorlinks=true,linkcolor=blue,citecolor=blue,urlcolor=blue]{hyperref}

\title{Combining C-LARA-2 and ChatGPT Voice Mode:\\Initial experiments\thanks{This report is based on and extends the paper \emph{Using ChatGPT Voice for spoken German practice: A case study}, submitted to EuroCALL 2026. At the time this report was posted, the authors did not know whether the EuroCALL paper would be published, following an apparent last-minute change in EuroCALL policy concerning whether an AI system may be listed as an author. The final submitted EuroCALL manuscript is preserved in the C-LARA-2 GitHub repository at \url{https://github.com/mannyrayner/C-LARA-2/tree/main/docs/publications/eurocall_2026_voice_mode}.}}
\author{%
  \parbox{0.96\textwidth}{\centering
    Sarah Wright\textsuperscript{a}, Manny Rayner\textsuperscript{b},
    ChatGPT C-LARA-Instance\textsuperscript{c}, and Branislav B\'{e}di\textsuperscript{d}\\[0.6em]
    \small \textsuperscript{a}Independent researcher, Australia,
    \href{mailto:sarah@tamargas.com.au}{sarah@tamargas.com.au}\\
    \small \textsuperscript{b}Independent researcher, Australia,
    \href{mailto:mannyrayner@c-lara.org}{mannyrayner@c-lara.org}\\
    \small \textsuperscript{c}OpenAI large language model instance participating in the C-LARA project,
    \href{mailto:chatgptclarainstance@proton.me}{chatgptclarainstance@proton.me}\\
    \small \textsuperscript{d}The \'{A}rni Magn\'{u}sson Institute for Icelandic Studies, Iceland,
    \href{mailto:branislav.bedi@arnastofnun.is}{branislav.bedi@arnastofnun.is}
  }%
}
\date{14 August 2026}

\begin{document}

\maketitle

\begin{abstract}
We would like C-LARA-2, an open-source platform for creating multimodal language-learning resources, to support free spoken conversation about its texts and images. This report presents an initial two-part pre-study. First, two anglophone learners used OpenAI's ChatGPT Voice Mode / Pro Tier for German conversation practice approximately once a week over eight weeks, usually in 60--90 minute sessions. Participant observations and post-session summaries suggest that the system supported confidence, fluency, situated vocabulary learning, flexible topic development, and low-pressure practice, while also revealing limitations in correction, speaker identification, learner tracking, factual reliability, and cost. Second, after improvements in ChatGPT's ability to access web pages removed an earlier practical obstacle, we tested a lightweight integration: a learner opens a C-LARA-2 text in one browser window, gives its URL to ChatGPT Voice Mode / Pro Tier in another, and then switches to spoken interaction. ChatGPT Voice Mode / Pro Tier can access the text, navigate its pages, view its images, and discuss the material. The consented session summaries are also stored in the C-LARA-2 repository, where the project Assistant can answer questions about them. The evidence remains anecdotal, but the combined experiments suggest that a usable spoken extension to multimodal C-LARA-2 materials may already be possible with modest interface engineering.
\end{abstract}

\noindent\textbf{Keywords:} spoken interaction; ChatGPT Voice Mode / Pro Tier; German language learning; foreign-language anxiety; conversational AI.

\section{Introduction and overview}

Learners often have only limited opportunities for spoken language practice in foreign and second language learning (L2), both in classroom environments and during self-study. Even learners who have studied an L2 for years and accumulated substantial passive knowledge may still, as a result of inadequate spoken practice, lack the confidence to use the language spontaneously. Many modern technologies support personalised and adaptive L2 learner experiences while providing opportunities for real-time interaction in both written and spoken form \citep{Urbaite2024}. In particular, technologies powered by generative artificial intelligence (GenAI) not only assist learners in meeting their overall learning goals in L2 but are also reported as contributing to increased speaking proficiency and reduced speaking anxiety \citep{DuDaniel2024,MengZou2026,SaarelaEtAl2026}.

The broader context for our work is C-LARA-2, a reimplementation of the open-source C-LARA platform for creating multimodal language-learning resources, including glossed texts, illustrated picture books, audio, and learner-specific resources \citep{BediEtAl2025}. We would like learners to be able not only to read and study these resources, but also to hold free spoken conversations about their texts and images. We therefore treat the work reported here as a pre-study for adding AI-based spoken interaction to C-LARA-2.

For the spoken-interaction component, we used ChatGPT Voice Mode / Pro Tier. This was partly a pragmatic choice: the authors already had access to the Pro Tier subscription for C-LARA-2 development, it offered direct real-time speech interaction without requiring us to build a separate speech pipeline, and preliminary trials with the less expensive Plus Tier had not supported the kind of extended free conversation we wanted to study. The relevant practical details and limitations are described below.

There are two obvious questions. First, how well does ChatGPT Voice Mode / Pro Tier work as a free conversation partner for language learning? Second, if its conversational performance is promising, how easily can it be combined with C-LARA-2's multimodal texts? When most of the pre-study was carried out, comparatively uninteresting technical limitations in ChatGPT's ability to access web pages made the second question difficult to investigate cleanly. We consequently focused first on conversational performance. More recent improvements appear to have removed the immediate access problem, allowing us to carry out an initial integration experiment as well.

The first part of the pre-study involved using ChatGPT Voice Mode / Pro Tier as a conversation partner to support L2 German speaking skills in unconstrained situations. The system we used was priced at USD~200 per month, and was primarily designed for coders and business users rather than as an educational language-learning product \citep{Reuters2024}. Although this cost represents a practical barrier to adoption, and could be regarded as contributing to the digital divide in language learning \citep{Yaman2015}, the platform is publicly available, which makes the study straightforward to replicate.

We should note two things at once. First, the case study is entirely anecdotal. It involves only one main subject, who needed to improve her spoken German conversational skills for practical reasons. She knew that one of the authors had access to ChatGPT Voice Mode / Pro Tier, and asked whether it would be possible to use it for these purposes; based on the remarkably positive results obtained even in the first session, we decided to keep records so that we could write up the outcomes. Second, it became clear during the study that ChatGPT Voice Mode / Pro Tier appeared in practice to be necessary for the type of free conversation practice described here. We experimented first with the much cheaper Plus Tier (USD~20 per month), and found that performance was clearly inadequate. The technology is developing rapidly, so both pricing and capability comparisons should be understood as time-specific observations rather than permanent conclusions.

Methodologically, the conversational case study is situated within the framework of individualised and self-regulated learning, which helps learners adapt to individual differences, manage their own learning processes, and gradually become more independent while achieving their learning goals \citep{Zimmerman2002}. Individual sessions using ChatGPT Voice Mode / Pro Tier create a unique microlearning environment for enhancing specific language skills \citep{KohnkeMoorhouse2026}. Building on this view, the study examines whether ChatGPT Voice Mode / Pro Tier can support speaking and listening practice during home self-study and partly take over the role of a human conversation tutor. The first part is explored through six questions: (1) whether ChatGPT Voice Mode / Pro Tier is sufficiently good at supporting realistic conversations in German; (2) whether it can adjust to the learners' level and maintain comprehensible conversation; (3) whether it can recognise what the learners say and respond accordingly; (4) how well it manages focus, follow-up questions, turn-taking, rephrasing, and learner-selected topics; (5) whether regular use can improve learners' confidence and fluency; and (6) what limitations remain. The second part asks whether ChatGPT Voice Mode / Pro Tier can access and discuss an existing C-LARA-2 text and whether the project's repository-grounded Assistant can use the resulting study records.

The rest of the report is organised as follows. Section~2 describes previous research in GenAI-supported spoken interaction in CALL. Section~3 describes the conversational study. Section~4 presents its findings. Section~5 reports the initial C-LARA-2 integration experiments. The final section summarises and concludes.

\section{Previous work}

Several recent papers suggest that combining spoken interaction and GenAI in CALL is a promising development that is starting to attract considerable interest. For example, \citet[p.~122]{Savvidou2026} claims that GenAI-powered CALL platforms like ChatGPT Voice Mode / Pro Tier, Duolingo Max, Gemini, Speak AI, Replica AI, and Talkpal AI can be helpful tools offering realistic speaking-practice opportunities, immediate feedback, personalised conversations, and opportunities for learners to continue practising independently outside the classroom, but unfortunately gives few details. In a study conducted by \citet{LuEtAl2026}, ChatGPT generates conversational scenarios while spoken responses from English language learners are captured by a speech-recognition system and the conversation is conducted with an animated conversational agent generated by another system called D-ID AI. Their results claim that repeated practice improves learners' speaking proficiency in the target language and reduces anxiety from having real-life conversations, but it is unclear from their paper what kind of spoken interaction is actually supported. The study we know of which is closest to ours is that of \citet{CarreraNunezEtAl2025}, which describes an experiment in which ChatGPT Voice Mode / Pro Tier was incorporated into L2 learning in regular and guided practice sessions in a classroom environment. Tests provided strong statistical evidence that it led to improvement of different language skills, particularly pronunciation, fluency, vocabulary, and grammar, as well as the reduction of language anxiety; the authors attribute this to the tool's non-judgemental environment.

Jumpspeak is a particularly relevant commercial comparison because it is explicitly designed around AI-supported spoken language practice \citep{Jumpspeak2026,JumpspeakAppStore2026}. Product descriptions and recent hands-on reviews suggest that its dialogue is genuinely generative: learners can select or create topics, speak rather than choose only fixed responses, and depart from suggested dialogue paths. At the same time, the interaction appears to remain organised primarily around role-play scenarios and language-learning exercises, often with scenario framing and suggested responses; reports also describe some user-created exchanges as short or basic and indicate limited continuity between sessions \citep{LangolyJumpspeak2026}. This suggests a useful hypothesis rather than a demonstrated ranking: Jumpspeak may support relatively free spoken production within an explicit AI-tutoring framework, whereas ChatGPT Voice Mode / Pro Tier may function more like a general conversational partner with language learning as a comparatively unobtrusive layer. Our observations support the ChatGPT side of this distinction, but we have not tested the systems head-to-head and should not claim superiority on untested dimensions. A controlled comparison of conversational freedom, continuity, correction, and learning outcomes would be valuable future work.\footnote{We considered a direct Jumpspeak comparison, but did not proceed after one author encountered an unpleasantly difficult subscription-cancellation process and decided not to resubscribe. Public consumer reviews contain similar allegations about cancellation and charges \citep{TrustpilotJumpspeak2026}; these user reports have not been independently verified and concern subscription management rather than the pedagogical quality of the system.}

All of the studies cited above stress how powerful the combination of GenAI and spoken interaction can be, particularly for learners with an autonomous learning style, and their overall conclusions agree well in principle with the findings we present here. The details, however, are less clear to us: crucially, we are uncertain about the extent to which the spoken conversations can reasonably be considered to be unconstrained, an issue which in practice is strongly related to the reliability of the spoken interaction. For obvious reasons, and particularly for maintaining student engagement, this is an important question; as already noted, our preliminary trials with lower-tier access gave clearly unsatisfactory results for unconstrained free conversation, while ChatGPT Voice Mode / Pro Tier performed much better. The technology concerned is evolving so quickly that this may only reflect a transient issue, but we note it for what it is worth.

\section{Background to the study}

As we said in the previous section, there is rapidly growing interest in the use of generative AI tools for speaking practice outside the classroom, but it is unclear to us that the most commonly available platforms can adequately support unconstrained conversation and function as true conversation partners during home self-study. Since our preliminary experiments suggested that this was not the case, the central research goal of the study reported here was to explore whether a high-end system, ChatGPT Voice Mode / Pro Tier, would be equal to the task. In particular, we wished to know whether it would allow learners to take part in spoken interaction, practise everyday situations, ask for vocabulary, and receive immediate responses.

The experiments used a small-scale qualitative case-study design and were carried out during eight regular sessions over approximately eight weeks, each lasting about 60--90 minutes. Participant~1, the main learner, was in her early twenties. She was a native speaker of English who had previously studied German for six years and been awarded a B1-level certificate reflecting reasonable low-intermediate skills in reading, writing, speaking, and listening; unfortunately, she had not maintained contact with the language for some time, and in particular felt that her fluency had declined considerably. She was, however, planning a trip where she would spend an extended period in German-speaking countries, and consequently wanted to improve her confidence in spoken German. For this reason, we agreed that we would begin by practising conversations organised around practical scenarios she could expect to encounter on her trip, and, if these went well, progress to more challenging topics.

Participant~2, in his mid-sixties, was a native speaker of English who had previously lived in several countries and worked in numerous jobs involving speech and language technology. He had near-native skills in Swedish, strong (C1) reading ability in French, Norwegian, and Danish, and B1/B2 reading ability in German, but his spoken German skills were only around A2 level. He participated as a facilitator and observer, helping to guide the conversation when needed; his main motivation for improving his German was to be able to watch German films more easily. The third participant, who acted as the external observer, was an experienced language teacher. He attended one early session and followed the remaining sessions through written summaries.

The sessions were not recorded or formally transcribed. After each session, ChatGPT generated a detailed written summary, which formed the main dataset. These summaries were analysed using directed qualitative content analysis \citep{AssarroudiEtAl2018}, which helped to extract relevant data based on predefined research criteria: the naturalness and real-life quality of the conversation, adaptation to the learner's language level, accuracy of speech recognition and appropriateness of responses, conversation management and learner engagement, and possible changes in the learners' confidence and fluency.

\section{Findings}

The sessions suggested several ways in which ChatGPT Voice Mode / Pro Tier can support spoken language practice. Since the study was only intended to be exploratory and offered anecdotal evidence, these benefits could not be measured quantitatively. They were, however, visible as the sessions unfolded and as the participants reflected on their learning experience. This section provides an interpretation of the findings divided into two categories: observed benefits and limitations.

\subsection{Observed benefits}

\subsubsection*{Authentic conversation quality leading to increased learner confidence}

The most striking benefit was Participant~1's increased willingness to speak. At the beginning of the experiment, her difficulty was not primarily a lack of German vocabulary or grammar. She had studied German at school and retained a reasonable passive understanding, but was hesitant about using the language spontaneously. ChatGPT Voice Mode / Pro Tier provided a setting in which she could speak without the social pressure of addressing an unfamiliar human interlocutor. It generally treated partially successful utterances as meaningful and responded to the intended content, thus minimising the stress associated with the fear of speaking incorrectly and, as a consequence, creating a low-pressure environment for fluency practice. After only two or three weeks, Participant~1 was already less dependent on rehearsed scenarios and more willing to initiate topics, ask questions, and continue speaking despite uncertainty.

The sessions provided a substantial amount of practice. Meeting approximately weekly for 60--90 minutes gave Participant~1 extended exposure to real-time spoken interaction. This is difficult to obtain in many ordinary learning contexts. A human conversation partner may not be available, and a teacher's time is usually limited. ChatGPT Voice Mode / Pro Tier could sustain long conversations without a fixed lesson plan and without the fatigue or scheduling constraints that often restrict spoken practice.

\subsubsection*{ChatGPT Voice Mode / Pro Tier adjusting to the learner's level by providing scaffolding}

A second important benefit was immediate scaffolding. When Participant~1 needed a word or phrase, ChatGPT Voice Mode / Pro Tier could usually supply it quickly and naturally within the conversation. Examples included practical vocabulary such as \emph{der Eiffelturm} for ``the Eiffel Tower'', \emph{das Armband} for ``bracelet'', \emph{die Aussprache} for ``pronunciation'', \emph{der Dieb} and \emph{der Diebstahl} for ``thief'' and ``theft'', \emph{das Katzenfutter} and \emph{das Katzenleckerli} for ``cat food'' and ``cat treats'', \emph{der Thunfisch} for ``tuna'', and \emph{die \'{U}berschwemmung} for ``flooding''. These items were not introduced as part of a pre-planned syllabus. They arose, rather, because the learner needed them in order to say something she spontaneously wanted to express. ChatGPT Voice Mode / Pro Tier also helped the learner remain in the target-language interaction when she lacked a precise form. In some cases the AI tool supplied a missing word; in others, the tool inferred the intended meaning and responded in simple German. This allowed the conversation to continue without repeatedly stopping for translation. The scaffolding was therefore not only lexical but also interactional.

\subsubsection*{Conversation management and learner engagement supporting role-play scenarios}

A third benefit was the system's flexibility. The early sessions rehearsed practical travel scenarios of the kind informally agreed at the beginning of the study, organised as role plays between the human and the AI (checking into a campsite, visiting a pharmacy, and so on). Later sessions became less scripted and more personal, talking about previous trips, food and drink, music, concerts, chess, fairy tales, films, pets, weather, pronunciation, and AI. This strongly suggests that ChatGPT Voice Mode / Pro Tier was not merely useful for practising set travel phrases, but could from the start support natural conversation. This flexibility is one of the distinctive strengths of a general conversational AI. A conventional coursebook can provide carefully graded scenarios, but it cannot easily follow the learner from a campsite role-play to a discussion of Gershwin, bad weather, cat treats, or \emph{2001: A Space Odyssey}. A human teacher can do this, but has limited time. ChatGPT Voice Mode / Pro Tier offered a form of abundant, responsive conversational practice that could adapt to the learner's interests in real time.

\subsection{Observed limitations}

\subsubsection*{Limited pedagogical control}

The most important limitation was the lack of explicit pedagogical control. ChatGPT Voice Mode / Pro Tier behaved primarily as a helpful conversation partner. This made the sessions relaxed and encouraging, but it also meant that the system did not consistently pursue language-learning goals. It did not maintain an explicit learner model, track Participant~1's recurring errors, or adapt its correction strategy over time. In particular, the AI usually prioritised communicative success over formal accuracy.

This observation is specific to the system available during the main conversational experiments. Since then, Manny Rayner has carried out an informal French test with ChatGPT version~5.6 and reports a marked improvement in pedagogically relevant behaviour, in particular the correction of errors involving grammatical gender and articles. This is an anecdotal observation rather than a controlled comparison, and it does not show that the other limitations have been resolved. It does, however, reinforce the importance of dating capability claims. Given the rapid improvements observed during the last two years, it is reasonable to expect further language-pedagogy capabilities to emerge as a by-product of general model improvements, without requiring each one to be engineered as a specialised CALL feature.

\subsubsection*{Speech recognition, clarification, and turn-taking}

A second set of limitations concerned the mechanics of spoken interaction. The speech input was sometimes noisy, incomplete, or incorrectly interpreted, especially when the speakers hesitated, code-switched, used proper names, or spoke in fragments. In several cases the AI responded to what it believed it had heard rather than to what the speaker had intended. Some misunderstandings were minor, but others temporarily diverted the conversation. Turn-taking also caused occasional problems. The system sometimes treated a pause as the end of a turn, produced a generic response where a more targeted follow-up would have been preferable, or decided to close the conversation prematurely with farewell formulas.

\subsubsection*{Confusing speaker voices}

The sessions usually involved Participant~1 and Participant~2 sitting together and both speaking to ChatGPT Voice Mode / Pro Tier. This created a further problem: the system did not reliably know who was speaking, which was sometimes an annoying limitation.

\subsubsection*{Factual reliability}

The conversations often moved beyond language practice into cultural, historical, linguistic, or technical information. In many cases the AI's responses were useful and plausible, but they were not systematically checked during the sessions. This creates a familiar risk: fluent conversational delivery can make uncertain information about diverse topics provided by ChatGPT Voice Mode / Pro Tier sound authoritative. In the opinion of the two participants, however, there were very few occasions where information provided by the AI was in fact dubious or incorrect.

\subsubsection*{Access and cost}

As noted, an extremely important practical limitation right now is cost, since acceptable performance for extended, flexible spoken interaction depends on having access to the expensive ChatGPT Voice Mode / Pro Tier. We only had a Pro Tier subscription because we use it for coding in the C-LARA project, and we think it unlikely that many people would purchase one just to use it as a conversation partner. In practice, we suspect that access will be limited to people who, like us, already have it for other reasons.

\section{Initial integration experiments with C-LARA-2}

The conversational results make it attractive to add a spoken-language component to C-LARA-2. The platform already supports multimodal pedagogical texts, while ChatGPT Voice Mode / Pro Tier supports sustained spoken interaction. A useful preliminary integration should let a learner open a C-LARA-2 text, discuss its words and images, and move naturally between reading and conversation without repeatedly copying material between systems.

\subsection{From screenshot transfer to URL access}

Our first experiment used two browser windows, one for C-LARA-2 and one for ChatGPT Voice Mode / Pro Tier. The learner opened a C-LARA-2 text page, took a screenshot, uploaded the screenshot to ChatGPT Voice Mode / Pro Tier, and then switched to spoken interaction. ChatGPT Voice Mode / Pro Tier could discuss both the text and the image, but the repeated screenshot workflow was far too cumbersome for practical use.

More recently, ChatGPT Voice Mode / Pro Tier's web-page access has improved enough to support a much cleaner experiment. The learner again opens the C-LARA-2 text and ChatGPT Voice Mode / Pro Tier in separate browser windows, but now only needs to give ChatGPT Voice Mode / Pro Tier the URL of the text before switching to Voice Mode. ChatGPT Voice Mode / Pro Tier can access the text directly, navigate between its pages, view its images, and use that material as the subject of the spoken conversation. No C-LARA-2 application code or dedicated interoperability layer is required for this preliminary workflow. In our initial tests it is already more or less usable, although a systematic evaluation across different texts, access settings, browsers, and conversation lengths remains to be carried out.

This is not yet a full product integration. The learner still has to coordinate two windows, provide the URL, and establish the desired task. A more polished interface should open or connect ChatGPT Voice Mode / Pro Tier from the C-LARA-2 text, keep page position and conversational context synchronized, make access failures visible, and offer explicit pedagogical controls. Nonetheless, the experiment changes the engineering assessment: the immediate problem no longer appears to be transferring every page manually, but designing a reliable and pedagogically useful interface around capabilities that are already available.

\subsection{Making the pre-study evidence accessible to the Assistant}

With Participant~1's consent, the eight post-session records used in this pre-study have been placed in the public C-LARA-2 GitHub repository. Seven are AI-generated retrospective summaries and one consists of contemporaneous human-authored notes; none is a recording or formal transcript. Repository storage makes the evidence inspectable by readers and also makes it available to the repository-grounded C-LARA-2 Assistant.

For example, we asked the Assistant:

\begin{quote}
\emph{In the experiments described in the EuroCALL 2026 Voice Mode paper, in which session does Sarah talk about visiting the Eiffel Tower? Can you find any details in the relevant session summary?}
\end{quote}

The Assistant identified Session~6, dated 14 May 2026, and linked its answer to the session index and relevant summary lines. It reported that Sarah described a previous Paris trip with her German exchange family, including a visit to the Eiffel Tower; extracted the vocabulary \emph{der Eiffelturm} and \emph{das Armband}; and noted her account of a bracelet scam that led into a discussion of tourist safety. Importantly, it also warned that the source was an AI-generated retrospective summary rather than a recording or transcript. The source records are available at \url{https://github.com/mannyrayner/C-LARA-2/tree/main/docs/publications/eurocall_2026_voice_mode/session_summaries}.

This example is modest, but it illustrates a second kind of integration. The conversational system can discuss a learner-facing C-LARA-2 text, while the project Assistant can retrieve and explain repository-held evidence about previous learning sessions. A future system could connect these functions more directly: post-session summaries could support vocabulary lists, corrected sentences, grammar explanations, follow-up exercises, and longitudinal learner models, subject to appropriate consent, privacy, and teacher oversight.

\subsection{Pedagogical and technical next steps}

A dedicated language-practice mode could allow learners or teachers to choose how often corrections are given, while delayed feedback could support fluency during the session and accuracy afterwards. The combined system should distinguish speakers, connect errors and vocabulary needs to the correct learner, and track progress across sessions. Conversations could be grounded in visible texts, images, and exercises while teachers retain control of learning goals, assessment, and final pedagogical decisions.

The two preliminary experiments therefore support a more concrete programme of work. Immediate next steps are to reproduce the URL workflow systematically, determine which published and access-controlled C-LARA-2 texts ChatGPT can reliably use, test navigation and image understanding, and study whether learners find two-window interaction acceptable. Longer-term work should replace the two-window arrangement with a controlled C-LARA-2 interface and evaluate learning outcomes with a larger and more varied participant group.

\section{Summary and conclusion}

We have presented a two-part pre-study exploring ChatGPT Voice Mode / Pro Tier as a German conversation partner and its preliminary integration with C-LARA-2 multimodal texts. We do not claim that the results are more than anecdotal. They refer to one main subject and are based on AI-generated post-session summaries and informal subject feedback rather than full recordings, transcripts, or independent speaking assessments, and cannot, therefore, demonstrate measurable improvements.

With the above caveats, we are much encouraged by our experience. ChatGPT Voice Mode / Pro Tier showed strong potential for low-pressure conversation practice. It was able to maintain extended and natural conversations about both constrained and open-ended topics. It could adapt its language to the learner's level by simplifying, paraphrasing, explaining, providing vocabulary, and generally reformulating when difficulties occurred. Speech recognition was mostly sufficient for the conversation to continue, although misunderstandings and requests for repetition still appeared, particularly when pronunciation or wording was unclear.

ChatGPT Voice Mode / Pro Tier also showed an ability to manage the conversation by asking follow-up questions, supporting turn-taking, maintaining selected topics, and elaborating its explanation when the learner did not understand. Its flexible role as conversation partner, vocabulary assistant, and provider of clarifications helped the learner to continue speaking without frequent interruption. Over the eight sessions, the learner became more willing to introduce topics, ask questions directly, and continue speaking despite limited vocabulary or grammatical uncertainty. She reported this as experiencing increased confidence and greater conversational fluency, and the observer concurred in this judgement. This is in line with \citet{CarreraNunezEtAl2025}, who find that regular sessions with GenAI Voice Mode over a longer period of time can contribute to increased language skills.

Although our contribution is modest, we interpret it as suggesting that ChatGPT Voice Mode / Pro Tier and similar high-end systems should be taken seriously as existing, already available providers of useful spoken language practice. They not only give instant access to a virtual conversation partner, imitate real-life scenarios, supply vocabulary on request, tolerate imperfect learner language, and allow discussion of personally meaningful topics, but also help reduce speaking anxiety; all of this agrees well with previous research \citep{DuDaniel2024,MengZou2026,SaarelaEtAl2026}. Since the platform we used is openly available, anyone can repeat our experiments. We would very much like to see other people doing so and reporting more substantial results. Most importantly, we wish to know whether they concur in our judgement that it can support long, unconstrained, pedagogically useful conversations with a non-trivial and varied sample of learners.

As we have said, although the mid-2026 version of ChatGPT Voice Mode / Pro Tier is already very useful as a conversation partner, it has clear limitations. It is not a specialised CALL tutor. It prioritises conversational flow over systematic correction to an extent that can be counterproductive, especially since it can misinterpret spoken input. It does not reliably distinguish speakers. Its pedagogical integration with text and images remains preliminary, although URL-based access now makes a practical two-window workflow possible. Above all, it is too expensive for most language learners to consider using it. One plausible interpretation of the current commercial landscape is that the strongest systems, including Jumpspeak, are constrained at least as much by the economics of providing sustained generative speech interaction as by fundamental technical limits; this hypothesis requires direct comparative evidence. There are nonetheless grounds for optimism: the cost of running models at a fixed level of capability has fallen sharply, with the 2025 AI Index reporting a more than 280-fold reduction in the inference cost of a system performing at approximately GPT-3.5 level between November 2022 and October 2024 \citep{MaslejEtAl2025}. If this price-performance trend continues alongside improvements in conversation, correction, context handling, and multimodal access, strong and affordable multimodal AI language tutors capable of free conversation may soon become widely available. The implications for spoken CALL would be considerable.

\section*{Division of work statement}

The role of Participant~1 (main learner) was occupied by Sarah Wright. The role of Participant~2 (co-learner) was occupied by Manny Rayner. The role of the External Observer was occupied by Branislav B\'{e}di. The role of the AI was occupied by ChatGPT C-LARA-Instance.

The first draft of the EuroCALL paper on which this report is based was written by ChatGPT C-LARA-Instance and then heavily revised, first by Branislav B\'{e}di, then by Manny Rayner, and last by Sarah Wright. The initial adaptation into the present report was prepared by OpenAI Codex from instructions supplied by Manny Rayner. The human authors, primarily Manny Rayner, have reviewed the report and take full responsibility for the version released publicly.

\section*{Funding}

The ChatGPT Voice Mode / Pro Tier subscription used for these experiments, and also for the development of C-LARA-2, is funded privately by Manny Rayner and Cathy Chua.

\bibliographystyle{plainnat}
\bibliography{references}

\end{document}
